\chapter{Conclusiones y trabajos futuros}\label{cap:conclusiones}

% Párrafo introductorio: dónde se llegó y qué queda por delante.
% Este capítulo tiene tanto peso como cualquier otro.

\section{Cumplimiento de los objetivos}

\Lucas{Redactar: contrastar lo logrado frente a los objetivos del capítulo 1 y los requisitos del capítulo 4 — qué entró en la prueba de concepto y qué quedó fuera.}

\section{Principales resultados y aportaciones}

\Lucas{Redactar: los resultados centrales del trabajo —herramienta con cinco comandos, cómputo y base de datos, AWS y Azure, demo de tres capas— y qué aporta frente al estado del arte (interoperabilidad real desde un manifiesto agnóstico).}

\section{Limitaciones del trabajo}

\Lucas{Redactar: limitaciones del alcance —dos proveedores, tipos de recurso acotados, sin almacenamiento de objetos ni análisis de seguridad avanzado—, asumidas conscientemente por el carácter de prueba de concepto.}

\section{Reflexión personal sobre el aprendizaje}

\Lucas{Redactar: reflexión personal sobre el proceso de aprendizaje —dificultades técnicas (diferencias AWS/Azure, referencias entre recursos), el trabajo con metodología iterativa y el uso de IA como apoyo—.}

\section{Trabajo futuro}

\Lucas{Redactar: líneas de evolución futura, alineadas con el capítulo 4 (GCP, interfaz web, componente agéntico, lock-in de código, monitorización, gestión de estado).}
